\section{The generator bound at an odd prime}
\label{sec:selmer-generators}

Let \(E\) be a number field of degree \(D\) admitting an embedding in
\(\mathbb Q(\mu_F)\), where \(F\ge1\). Fix a prime \(\ell\) and let
\(S_E\) be the finite places of \(E\) above \(\ell\). Write
\(G_{E,S_E}\) for the Galois group unramified at all finite places
outside \(S_E\). Take the discrete module
\(A=(\mathbb Q_\ell/\mathbb Z_\ell)(-k)\) for an integer \(k\).
We prove that the subgroup
\[
        V=H^1(G_{E,S_E},A)[\ell]
\]
is finite of \(\mathbb Z_\ell\)-length \(g\), with
\[
              g\le 8\ell D\log(\ell F).
\]
This is the ambient cohomology group in
\eqref{eq:galois-annihilator-length}. In particular the estimate
allows ramification at places above \(\ell\) and imposes no strict
local condition at those places. The same calculation works at
\(\ell=2\) for this generator bound.

\subsection{Restriction to the first cyclotomic layer}

Put \(K=E(\mu_\ell)\), \(n=[K:\mathbb Q]\), and
\(S=\{\mathfrak p\mid\ell\text{ in }K\}\). The action on the
\(\ell\)-th roots embeds \(\operatorname{Gal}(K/E)\) in
\(\mathbb F_\ell^\times\), so \([K:E]\mid\ell-1\) and
\(n\le\ell D\). Both \(E\) and the required roots embed compatibly
in \(\mathbb Q(\mu_{\ell F})\), by extending the given embedding of
\(E\) to an algebraic closure and taking the compositum there.

For each \(v\in V\), choose a continuous representative \(z\).
The equation \(\ell v=0\) gives \(\ell z=db\), with coboundary
\(db(\sigma)=\sigma b-b\). Multiplication by \(\ell\) on \(A\)
is onto. Choosing \(b'\) with \(\ell b'=b\), we may replace \(z\)
by \(z-db'\), so its values belong to \(A[\ell]\).
Fix an additive identification of this group with \(\mathbb F_\ell\).
On \(G_K=\operatorname{Gal}(\overline E/K)\) the residual
cyclotomic character is trivial. The restriction of the inflated
cocycle is therefore an additive continuous homomorphism \(t_v\).
It vanishes on inertia outside \(S\), since such inertia maps into inertia
at a place outside \(S_E\), which the original quotient kills.

The choices of representatives give an injection of sets from \(V\)
to these homomorphisms. Indeed, if the restrictions for two classes
agree, their difference \(f\) vanishes on \(G_K\) and has values in
\(A[\ell]\). Its cocycle law makes it constant on right cosets of
\(G_K\), and the action on its values factors through
\(\Delta=\operatorname{Gal}(K/E)\). Write \(d=|\Delta|\) and
\(b=-d^{-1}\sum_{\tau\in\Delta} f(\tau)\) in \(A[\ell]\), where
\(d\) is prime to \(\ell\). Summing
\(f(\sigma\tau)=f(\sigma)+\sigma f(\tau)\) over \(\tau\) gives
\(f(\sigma)=\sigma b-b\). The two original classes coincide.
This proof uses the homomorphisms only as an upper bound on the
number of classes and is independent of the choices made for
representatives.

\subsection{Counting the permitted Kummer classes}

Choose a primitive root \(\zeta_\ell\in K\). For each homomorphism
\(t:G_K\to\mathbb F_\ell\) just described, Hilbert's theorem 90
on its finite cyclic quotient gives \(\beta\ne0\) in the field
fixed by \(\ker t\) such that
\(\sigma\beta/\beta=\zeta_\ell^{t(\sigma)}\).
Indeed these roots form a multiplicative cocycle since
\(G_K\) fixes \(\zeta_\ell\). Then \(a=\beta^\ell\in K^\times\).
Changing \(\beta\) by a scalar in \(K^\times\) changes \(a\) by
an \(\ell\)-th power. Conversely, if \(a'/a=c^\ell\) with
\(c\in K^\times\), then \(\beta'/(\beta c)\in\mu_\ell\subset K\),
and the two homomorphisms coincide. For \(\mathfrak p\notin S\)
the field fixed by \(\ker t\) is unramified at \(\mathfrak p\).
At any place above it the normalized integer valuation therefore
restricts to \(v_{\mathfrak p}\). It follows that
\(v_{\mathfrak p}(a)=\ell v(\beta)\) is divisible by \(\ell\).
Thus we have injected the homomorphisms into the finite group
\[
 \mathcal K_S=\{a\in K^\times:
          \ell\mid v_{\mathfrak p}(a)\text{ for }\mathfrak p\notin S\}
                      /K^{\times\ell}.
\]

Let \(U_S\) be the \(S\)-unit group and \(\mathrm{Cl}_S\) the
ideal class group of the ring of \(S\)-integers. Factoring
\((a)=\mathfrak a^\ell\) in that ring gives an exact sequence
\[
 1\longrightarrow U_S/U_S^\ell\longrightarrow \mathcal K_S
       \longrightarrow\mathrm{Cl}_S[\ell]\longrightarrow 1 .
\]
Indeed unique ideal factorization determines \(\mathfrak a\).
Multiplication of \(a\) by \(c^\ell\) changes \(\mathfrak a\) by
\((c)\). A class mapping to the identity can thus be represented by
an \(S\)-unit. An \(\ell\)-th root in \(K^\times\) of an \(S\)-unit
is itself an \(S\)-unit by its valuations. Finally any ideal class
of order dividing \(\ell\) has a representative whose \(\ell\)-th
power is principal, proving surjectivity.

If \(K\) has signature \((r_1,r_2)\), Dirichlet's \(S\)-unit theorem
gives free rank \(r_1+r_2-1+|S|\). The torsion is cyclic and contributes
at most \(\ell\) modulo \(\ell\)-th powers. We have \(r_1+r_2\le n\)
and \(|S|\le n\), the latter by the prime decomposition of \(\ell\).
Also \(\mathrm{Cl}_S\) is a quotient of the ordinary ideal class
group, of order \(h_K\). Hence
\[
      |V|\le|\mathcal K_S|\le\ell^{2n} h_K .
\]

\subsection{The class number and the uniform constant}

We give the estimate on \(h_K\) used here. For any degree-\(n\)
number field of discriminant of absolute value \(\Delta_K\), every
ideal class has an integral representative of norm
\(\le X=\sqrt{\Delta_K}\). Indeed Minkowski's bound supplies the
factor \((4/\pi)^{r_2}n!/n^n\le1\). To verify the last inequality,
\(r_2\le n-1\), and \(2^{n-1}n!\le n^n\) by induction from \(n=1\),
using \((1+1/n)^n\ge2\).

For each rational prime \(p\), fix an ordering of the primes above
it with residue degrees \(f_{p,i}\), and pad the list to length \(n\)
with empty positions. Encode an integral ideal with exponents
\(a_{p,i}\) as an ordered \(n\)-tuple of positive integers. Its
\(i\)-th entry has \(p\)-adic valuation \(f_{p,i}a_{p,i}\) in an
occupied position and zero in an empty one. Prime factorization
recovers the ideal, and the tuple product equals its norm. The number
of tuples with product at most \(X\) is bounded by
\[
 X\left(\sum_{1\le j\le X}\frac1j\right)^{n-1}
       \le X(1+\log X)^{n-1}.
\]
For the first inequality, fix the first \(n-1\) entries and bound
the last by \(X\) divided by their product. For the second, compare
the harmonic sum after its first term with an integral.
Since \(X\ge1\), taking logarithms and increasing \(n-1\) to \(n\)
gives
\[
  \log h_K\le \tfrac12\log\Delta_K+
                       n\log(1+\tfrac12\log\Delta_K).
\]

Now \(K\subseteq\mathbb Q(\mu_M)\) for \(M=\ell F\ge2\).
The conductor--discriminant formula expresses \(\Delta_K\) as the
product of the conductors of the \(n\) characters of
\(\operatorname{Gal}(K/\mathbb Q)\). Each divides \(M\), so
\(\Delta_K\le M^n\). Also \(n\le\varphi(M)\le M\). Since
\(\log M\le M-1\),
\[
   1+\tfrac12\log\Delta_K\le 1+\tfrac12 M(M-1)\le M^2.
\]
The preceding class-number inequality yields
\(\log h_K\le(5/2)n\log M\le3n\log M\).
In particular the bound is uniform in the subfield and the local data.

The group \(V\) is a finite vector space over \(\mathbb F_\ell\).
Its dimension is its \(\mathbb Z_\ell\)-length \(g\), so
\[
 g\log\ell=\log|V|\le2n\log\ell+\log h_K
          \le 5n\log(\ell F).
\]
Convexity of \(u\mapsto1/u\) on \([1,2]\) and its tangent at \(3/2\)
give \(\log2\ge2/3>5/8\). Division by \(\log\ell\ge\log2\) and
\(n\le\ell D\) prove the stated estimate. Its proof applies before
restricting to the eigenspace or imposing the decomposition-group
conditions defining the strict subgroup.
\footnote{\href{https://github.com/Deicyde/riemann-hypothesis/blob/40e56bab026552e089d7afcfeff08d2defb9c005/src/Development/AutomorphicArithmetic/AutomorphicVarieties/OddResidueBounds/Imports.lean}{Assembly of the generator estimate},
\href{https://github.com/Deicyde/riemann-hypothesis/blob/40e56bab026552e089d7afcfeff08d2defb9c005/src/Development/AutomorphicArithmetic/AutomorphicVarieties/OddResidueBounds/UnrInj.lean}{restriction to the cyclotomic layer},
\href{https://github.com/Deicyde/riemann-hypothesis/blob/40e56bab026552e089d7afcfeff08d2defb9c005/src/Development/AutomorphicArithmetic/AutomorphicVarieties/OddResidueBounds/UnrSel.lean}{Kummer injection},
\href{https://github.com/Deicyde/riemann-hypothesis/blob/40e56bab026552e089d7afcfeff08d2defb9c005/src/Development/AutomorphicArithmetic/AutomorphicVarieties/OddResidueBounds/SelCard.lean}{the local conditions and the unit bound},
\href{https://github.com/Deicyde/riemann-hypothesis/blob/40e56bab026552e089d7afcfeff08d2defb9c005/src/Development/AutomorphicArithmetic/CyclotomicBounds/ClassNumberBounds/IdealCounting/ClassC.lean}{Minkowski and ideal counting},
\href{https://github.com/Deicyde/riemann-hypothesis/blob/40e56bab026552e089d7afcfeff08d2defb9c005/src/Development/AutomorphicArithmetic/AutomorphicVarieties/OddResidueBounds/LogH.lean}{the cyclotomic class-number estimate}.}
